\documentclass[11pt]{article}
\usepackage[a4paper,margin=27mm]{geometry}
\usepackage{amsmath,amssymb,booktabs,microtype}
\usepackage[hidelinks]{hyperref}
\usepackage[numbers,sort&compress]{natbib}

\hypersetup{
  pdftitle={A computer-assisted three-state counterexample to a nonlinear phase-lag bound},
  pdfauthor={Xiaojian Zeng}
}

\newcommand{\dd}{\mathrm{d}}
\newcommand{\ee}{\mathrm{e}}
\newcommand{\norm}[1]{\lVert#1\rVert}

\title{A computer-assisted three-state counterexample to a nonlinear phase-lag bound}
\author{Xiaojian Zeng\\
\small PBC School of Finance, Tsinghua University, Beijing 100190, China\\
\small \texttt{xiaojianzeng@foxmail.com}}
\date{}

\begin{document}
\maketitle

\begin{abstract}
A phase-lag bound was proved for periodically driven two-state Markov systems and
conjectured for arbitrary finite networks by replacing the single rate prefactor with
the smallest edge prefactor.  We disprove the conjecture already on the complete
three-state graph.  The counterexample retains local detailed balance, all three
edges, time-independent positive prefactors, and a two-harmonic energy drive.  A
computer-assisted residual argument encloses the exact periodic state uniformly in
time.  Outward-rounded ball arithmetic certifies that both phase integrals are
positive while the proposed bound margin is strictly negative.  Thus the phase
exceeds the proposed upper bound.  The result shows that the slowest bare edge
prefactor alone does not control nonlinear multistate phase lag; no replacement bound
is claimed.
\end{abstract}

\noindent\textbf{Keywords:} Stochastic thermodynamics; Dynamical processes; Exact
results; Stationary states

\section{Introduction}

Periodic driving produces an intrinsic lag between a Markov system and its changing
energy landscape.  Deg\"unther, Koyuk and Seifert introduced a scalar phase for this
lag and proved a nonlinear upper bound for two-state systems
\cite{deguenther2022phase}.  For a general network with heterogeneous edge
prefactors, they proposed the extension
\begin{equation}
  \Delta\varphi\leq \arctan\!\left(\frac{\Omega}{k_{0,\min}}\right),
  \label{eq:conjecture}
\end{equation}
supported by random three-state calculations.  The same statement was still described
as a conjecture in the later thesis of Deg\"unther
\cite{deguenther2024bounds}.  The two-state proof does not directly extend because a
multistate network has several energy differences, rate splittings, and relaxation
channels.

Here we give one fully specified complete-three-state system for which
Eq.~\eqref{eq:conjecture} fails.  A floating-point solve is used only to select a
comparison path.  The mathematical conclusion follows from interval enclosures of
its continuum residual, an analytic contraction estimate for the exact period map,
and interval evaluation of the observables.  This separates discovery from the
certificate and rules out an apparent violation caused by time discretization or an
imperfectly periodic trajectory.

\section{Model and the proposed bound}

Consider states \(i\in\{1,2,3\}\), with all three undirected edges present and
\(k_{\mathrm B}T=1\).  For \(i<j\), let
\begin{align}
 k_{ij}(t)&=k_{0,ij}\exp\!\left[\alpha_{ij}(E_i(t)-E_j(t))\right],
 \label{eq:rate-forward}\\
 k_{ji}(t)&=k_{0,ij}\exp\!\left[-(1-\alpha_{ij})(E_i(t)-E_j(t))\right],
 \label{eq:rate-backward}
\end{align}
where \(k_{0,ij}>0\) and \(0\leq\alpha_{ij}\leq1\) are constant.  Hence
\(k_{ij}/k_{ji}=\exp(E_i-E_j)\).  The energies have a finite Fourier expansion,
\begin{equation}
 E_i(t)=E_i^0+\sum_{n=1}^{M}
 \left[c_{ni}\sin(n\Omega t)+d_{ni}\cos(n\Omega t)\right].
 \label{eq:drive}
\end{equation}
The probability column \(p\) obeys
\begin{equation}
 \dot p_i=\sum_{j\ne i}(p_jk_{ji}-p_ik_{ij}),\qquad
 p(t+T)=p(t),\quad T=\frac{2\pi}{\Omega}.
 \label{eq:master}
\end{equation}
Strictly positive continuous rates make the one-period stochastic matrix positive;
therefore its normalized fixed point exists, is unique, and is componentwise positive.

Define the period averages
\begin{equation}
 A=\frac1T\int_0^T p(t)\mathbin{\cdot}\ddot E(t)\,\dd t,
 \qquad
 B=\frac1T\int_0^T p(t)\mathbin{\cdot}\dot E(t)\,\dd t.
 \label{eq:AB}
\end{equation}
On the nondegenerate branch \(A>0,B>0\), the multistate phase of
Ref.~\cite{deguenther2022phase} is
\begin{equation}
 \Delta\varphi=\arctan\!\left(\frac{\Omega B}{A}\right).
 \label{eq:phase}
\end{equation}
Because \(\arctan\) is increasing and \(k_{0,\min}>0\), Eq.~\eqref{eq:conjecture}
is then equivalent to
\begin{equation}
 A-k_{0,\min}B\geq0.
 \label{eq:margin}
\end{equation}

\section{Exact counterexample}

Set \(M=2\), \(\Omega=6\), and use the exact terminating decimals
\begin{equation}
\begin{array}{c|ccc}
\text{edge }ij & 12&13&23\\ \hline
k_{0,ij}&1&300&8000\\
\alpha_{ij}&0.01&0.83&0.93
\end{array}
\label{eq:edgeparameters}
\end{equation}
with \(\theta=\Omega t\) and
\begin{align}
E_1(\theta)={}&-13+0.8\sin\theta+1.1\sin2\theta
                 -1.7\cos\theta+6.8\cos2\theta,
                 \label{eq:E1}\\
E_2(\theta)={}&-23+0.85\sin\theta+6.1\sin2\theta
                 +1.7\cos\theta-3\cos2\theta,
                 \label{eq:E2}\\
E_3(\theta)={}&0.
\label{eq:E3}
\end{align}
Every edge is retained, every prefactor is fixed and positive, and every splitting
parameter lies strictly inside its allowed interval.  In particular,
\(k_{0,\min}=1\).

Our central result is the following.

\medskip
\noindent\textbf{Proposition.}
\emph{The exact periodic state of Eqs.~\eqref{eq:rate-forward}--\eqref{eq:E3}
satisfies \(A>0\), \(B>0\), and \(A-B<0\).  Consequently it violates
Eq.~\eqref{eq:conjecture}.}
\medskip

The proof is computer-assisted; all discretization and rounding errors used in the
conclusion are enclosed below.

\section{Residual certificate}

\subsection{Periodic comparison path}

In the phase variable, Eq.~\eqref{eq:master} is
\begin{equation}
 p'(\theta)=G(\theta)p(\theta),\qquad G(\theta)=L(\theta)/\Omega,
 \label{eq:thetaode}
\end{equation}
where the columns of \(L\) sum to zero.  A numerical Radau solve supplies values and
slopes on 2048 uniform segments.  The final value and slope are set equal to the
initial ones, and cubic Hermite interpolation defines \(\widehat p_1,\widehat p_2\);
we set \(\widehat p_3=1-\widehat p_1-\widehat p_2\).  Thus the comparison path is
exactly normalized, periodic, and \(C^1\), regardless of the accuracy of the
proposal solve.  Each binary floating-point node and slope is converted to its exact
integer ratio before evaluation in Arb ball arithmetic
\cite{johansson2017arb}.

Let
\begin{equation}
 r(\theta)=\widehat p'(\theta)-G(\theta)\widehat p(\theta).
\end{equation}
Each Hermite segment is divided into 1024 closed subintervals.  At 256-bit working
precision, outward-rounded interval evaluation on all \(2{,}097{,}152\) boxes gives
\begin{equation}
 R:=\int_0^{2\pi}\norm{r(\theta)}_1\,\dd\theta
 \leq 8.40694407275998\times10^{-5}.
 \label{eq:R}
\end{equation}

\subsection{From residual to the exact periodic state}

For a continuous-time Markov generator with directed rates \(q_{ij}\), the standard
coupling estimate contracts zero-sum signed measures in \(\ell^1\) at the instantaneous
rate
\begin{equation}
 \gamma=\min_{i\ne j}\left(q_{ij}+q_{ji}
 +\sum_{\ell\ne i,j}\min(q_{i\ell},q_{j\ell})\right).
 \label{eq:dobrushin}
\end{equation}
Indeed, for a pair of starting states \(i,j\), the displayed rate is the sum of
direct jumps between the pair and the maximum common jump rate into every third
state; an infinitesimal coupling followed by integration gives the contraction.
For the three unordered pairs, the bracket is bounded below respectively by
\(k_{12}\), \(k_{31}\), and \(k_{32}\) in physical time.  Direct triangle bounds on
Eqs.~\eqref{eq:E1} and \eqref{eq:E2} give
\begin{equation}
 E_1-E_2\geq-8.25,\qquad E_1\leq-2.6,
 \qquad E_2\leq-11.35.
\end{equation}
It follows from Eqs.~\eqref{eq:rate-forward} and \eqref{eq:rate-backward} that
\(k_{12}\geq\exp(-33/400)\), whereas \(k_{31}>1\) and \(k_{32}>1\).
The exact one-period map therefore contracts the zero-sum subspace by at most
\begin{equation}
 q=\exp\!\left[-\frac{2\pi}{6}\ee^{-33/400}\right]
 \leq0.381260853664531<1.
 \label{eq:q}
\end{equation}

Put \(e=p-\widehat p\).  Variation of constants, periodicity of both paths, and
Eq.~\eqref{eq:q} imply
\begin{equation}
 \norm{e(0)}_1\leq\frac{R}{1-q},\qquad
 \sup_{0\leq\theta\leq2\pi}\norm{e(\theta)}_1
 \leq\frac{R}{1-q}+R
 \leq2.19941627295248\times10^{-4}.
 \label{eq:stateerror}
\end{equation}
This step connects the finite comparison path to the unique exact periodic state; no
numerical monodromy matrix or numerical ODE tolerance is assumed to be rigorous.

\subsection{Observable enclosure}

The Hermite cubics are integrated in closed form against the two Fourier harmonics,
with all evaluations enclosed by Arb.  For the comparison path this gives
\begin{equation}
 \widehat A=0.4997274966496502\ldots,\qquad
 \widehat B=6.1308368530420400\ldots .
\end{equation}
In \(\theta\) units the exact coefficients give
\begin{equation}
 \max_i\norm{\partial_\theta E_i}_\infty\leq20.75,
 \qquad
 \max_i\norm{\partial_\theta^2 E_i}_\infty\leq38.95.
\end{equation}
Combining these bounds with Eq.~\eqref{eq:stateerror} yields
\begin{align}
 |A-\widehat A|&\leq\Omega^2(38.95)\sup_\theta\norm{e(\theta)}_1
 \leq0.308402149793396,\\
 |B-\widehat B|&\leq\Omega(20.75)\sup_\theta\norm{e(\theta)}_1
 \leq0.027382732598259.
\end{align}
The final outward-rounded enclosures are
\begin{align}
 A&\in[0.1913253462459419,\;0.8081296470533585],\\
 B&\in[6.1034541204110998,\;6.1582195856729802],\\
 A-B&\in[-5.966894239630766,\;-5.295324473154014].
 \label{eq:finalintervals}
\end{align}
The first two lower endpoints establish the stipulated positive branch; the negative
upper endpoint in the third line proves the proposition.  A separate run with 4096
Hermite segments, 512 residual subboxes per segment, and 320-bit precision again
certifies all three strict signs.

\section{Physical meaning and limitations}

The counterexample is not produced by deleting a slow edge or by changing the phase
convention.  It is a complete, irreducible local-detailed-balance network on the same
positive branch for which the bound was proposed.  The witness combines heterogeneous
prefactors, asymmetric rate splittings, and two drive harmonics; the certificate does
not identify any one of these features as necessary.  What it establishes is narrow
but decisive: the slowest bare prefactor is insufficient as a universal nonlinear
multistate phase control.

We do not classify the set of violating drives, propose an alternative network
timescale, or address a different phase observable.  The example also does not
invalidate the proven two-state results of Ref.~\cite{deguenther2022phase}.  Finding a
replacement bound may require information from the full generator, its time-dependent
contraction, or network topology rather than one edge prefactor.

\section*{Data and code availability}

The exact-decimal parameters, the certificate program, the primary 256-bit output, a
changed-partition 320-bit replay, and environment instructions accompany this draft.
The certificate program regenerates the comparison path and performs every interval
evaluation used in Eqs.~\eqref{eq:R}--\eqref{eq:finalintervals}.

\section*{Acknowledgements and AI-use disclosure}

OpenAI Codex using GPT-5.6 Sol (OpenAI) assisted with literature support, scientific
reasoning, code, interval computation, an independent AI audit, and manuscript
drafting.  It is not an author.  The named author must critically review the scientific
content, references, code, disclosure, and final wording before any submission.  No
external submission has been made as part of this work.

\bibliographystyle{unsrtnat}
\bibliography{references}

\end{document}
